\documentclass[a4paper, serif, 10pt]{moderncv}

%% moderncv
% style and color
\moderncvstyle{banking}
\moderncvcolor{black}

% renew moderncv command to adapt for CJK colon
\renewcommand*{\cvitem}[3][.25em]{%
  \ifstrempty{#2}{}{\hintstyle{#2}：}{#3}%
  \par\addvspace{#1}}

\renewcommand*{\cvitemwithcomment}[4][.25em]{%
  \savebox{\cvitemwithcommentmainbox}{\ifstrempty{#2}{}{\hintstyle{#2}：}#3}%
  \setlength{\cvitemwithcommentmainlength}{\widthof{\usebox{\cvitemwithcommentmainbox}}}%
  \setlength{\cvitemwithcommentcommentlength}{\maincolumnwidth-\separatorcolumnwidth-\cvitemwithcommentmainlength}%
  \begin{minipage}[t]{\cvitemwithcommentmainlength}\usebox{\cvitemwithcommentmainbox}\end{minipage}%
  \hfill% fill of \separatorcolumnwidth
  \begin{minipage}[t]{\cvitemwithcommentcommentlength}\raggedleft\small\itshape#4\end{minipage}%
  \par\addvspace{#1}}

%% page layout/margins
\usepackage[top=2.5cm, bottom=2.5cm, left=1.5cm, right=1.5cm]{geometry}

\nopagenumbers{}

%% Babel config for English language
\usepackage[english]{babel}

%% fontspec
\usepackage{fontspec}

\IfFontExistsTF{Linux Libertine}{
  \setmainfont[Ligatures={TeX, Common}, Numbers=Lining, ItalicFont=Linux Libertine]{Linux Libertine}
}{}
\IfFontExistsTF{Linux Libertine O}{
  \setmainfont[Ligatures={TeX, Common}, Numbers=Lining, ItalicFont=Linux Libertine O]{Linux Libertine O}
}{}

%% CTeX
% CJK support, used to show CJK characters in the resume
%
% - fontset=none: disable builtin fontset but instead set the CJK font manually
% - heading=false: disable ctex heading
% - punct=kaiming: use kaiming punctuations styles for CJK
% - scheme=plain: use plain scheme, do not override \normalsize font size
% - space=auto: space settings for CJK characters
%
% ref:
% - http://ctan.mirrorcatalogs.com/language/chinese/ctex/ctex.pdf
\usepackage[UTF8, heading=false, punct=kaiming, scheme=plain, space=auto]{ctex}

\IfFontExistsTF{Noto Serif CJK SC}{
  \setCJKmainfont{Noto Serif CJK SC}
}{}
\IfFontExistsTF{Noto Sans CJK SC}{
  \setCJKsansfont{Noto Sans CJK SC}
}{}

%% line spacing
\usepackage{setspace}
\setstretch{1.125}

%% URL styling - use normal text instead of monospace
\urlstyle{same}

% Auto-underline all links
% Defer until \begin{document} so hyperref is already loaded
\AtBeginDocument{%
  \let\oldhref\href
  \renewcommand{\href}[2]{\oldhref{#1}{\underline{#2}}}%
}

%% Patch \httplink and \httpslink to handle full URLs
% The original moderncv macros always prepend http:// or https:// to URLs.
% This causes issues when the URL already has a protocol (e.g., https://example.com)
% resulting in malformed URLs like https://https://example.com.
% This patch checks if the URL already contains "://" and uses it directly if so.
% Note: We use \str_set:Nx (x-expansion) to expand macros like \@homepage first.
\ExplSyntaxOn
\str_new:N \g_yamlresume_url_str

\RenewDocumentCommand{\httpslink}{O{}m}{
  \str_set:Nx \g_yamlresume_url_str {#2}
  \str_if_in:NnTF \g_yamlresume_url_str {://}
    { \url{#2} }
    { \url{https://#2} }
}

\RenewDocumentCommand{\httplink}{O{}m}{
  \str_set:Nx \g_yamlresume_url_str {#2}
  \str_if_in:NnTF \g_yamlresume_url_str {://}
    { \url{#2} }
    { \url{http://#2} }
}
\ExplSyntaxOff

\name{陳昱翰}{}
\title{資深專案經理 | 敏捷開發與跨部門協作}
\phone[mobile]{+886 912 345 678}
\email{yuhan.chen@example.com}
\homepage{https://www.yuhanchen-pm.com}

\address{中國臺灣，台北市，台北市，台北市信義區松高路 1 號，110}{}{}

\extrainfo{{\small \faLinkedin}\ \href{https://www.linkedin.com/in/yuhan-chen-pm}{@yuhan-chen-pm} {} {} {} • {} {} {} 
{\small \faGithub}\ \href{https://github.com/yuhanchen-pm}{@yuhanchen-pm}}

\begin{document}

\maketitle

\section{簡介}

\cvline{}{\begin{itemize}
\item 擁有 8 年以上軟體與數位產品專案管理經驗，擅長從需求訪談、範疇定義到交付驗收的完整生命週期。
\item 精通 Agile/Scrum 與 Waterfall 方法論，能依專案特性調整整體節奏與資源配置。
\item 具備跨部門溝通與利害關係人管理能力，成功協調工程、設計、行銷與營運團隊達成目標。
\item 數據驅動決策，熟悉 Jira、Confluence、Asana 與 Tableau，持續優化流程與交付品質。
\end{itemize}}
\section{教育背景}

\cventry{2014年9月–2016年6月}
        {碩士，資訊管理學系，成績：4.0}
        {國立臺灣大學}
        {\url{https://www.ntu.edu.tw}}
        {}
        {\begin{itemize}
\item 碩士論文以「敏捷式專案管理在台灣軟體產業之應用」為題，訪談 12 家企業。
\item 擔任資訊管理研究所學會會長，籌辦產學講座與企業參訪。
\item 獲選為管理學院優秀研究生，並取得斐陶斐榮譽會員資格。
\end{itemize}
\textbf{課程}：專案管理、資訊系統分析與設計、資料庫管理、企業流程再造、敏捷軟體開發}
\section{工作經歷}

\cventry{2021年3月至今}
        {資深專案經理}
        {智邦科技股份有限公司}
        {\url{https://www.accton.com.tw}}
        {}
        {\begin{itemize}
\item 負責企業級網通軟體專案，管理 5 個跨功能團隊、共 40 位成員，年度專案預算達新台幣 8,000 萬元。
\item 導入 Scrum 與 Kanban 混合式管理，將版本交付週期由 12 週縮短至 8 週，提升 33\% 交付效率。
\item 建立標準化風險管理流程，提前識別關鍵風險並制定緩解計畫，降低 25\% 專案延遲率。
\item 主導與歐美客戶的季度業務審查，協調需求優先級與範疇變更，維持客戶滿意度 95\% 以上。
\item 輔導 3 位初階專案經理取得 PMP 證照，建立內部專案管理知識庫。
\end{itemize}
\textbf{關鍵字}：敏捷開發、跨部門協作、風險管理、利害關係人溝通、預算控管}

\cventry{2017年7月–2021年2月}
        {專案經理}
        {趨勢科技股份有限公司}
        {\url{https://www.trendmicro.com}}
        {}
        {\begin{itemize}
\item 管理雲端安全產品專案，從需求確認、開發排程到全球發布，年均完成 6 個主要版本。
\item 與產品經理合作定義產品路線圖，將客戶回饋轉化為可執行的使用者故事與驗收條件。
\item 協調 QA、DevOps 與技術支援團隊，建立自動化測試與持續整合流程，減少 40\% 回歸測試時間。
\item 每月向高階管理層彙報專案健康度、資源使用與里程碑達成率。
\end{itemize}
\textbf{關鍵字}：雲端安全、產品交付、持續整合、資源規劃}

\cventry{2016年7月–2017年6月}
        {專案管理師}
        {資拓宏宇國際股份有限公司}
        {\url{https://www.iisigroup.com}}
        {}
        {\begin{itemize}
\item 參與政府資訊系統標案，負責需求訪談、規格書撰寫與專案進度追蹤。
\item 協助導入 Jira 與 Confluence，統一專案文件與問題追蹤流程。
\item 支援驗收測試與教育訓練，確保系統順利上線並移交維運團隊。
\end{itemize}
\textbf{關鍵字}：政府標案、需求分析、文件管理、驗收測試}
\section{語言}

\cvline{中文}{母語或雙語水平 \hfill \textbf{關鍵字}：母語}
\cvline{英語}{完全專業水平 \hfill \textbf{關鍵字}：TOEIC 950、商業談判、跨國團隊協作}
\cvline{日語}{有限工作水平 \hfill \textbf{關鍵字}：JLPT N3、日常會話}
\section{專業技能}

\cvline{專案管理}{專家 \hfill \textbf{關鍵字}：Agile、Scrum、Waterfall、風險管理、範疇管理、預算控管}
\cvline{工具與平台}{高級 \hfill \textbf{關鍵字}：Jira、Confluence、Asana、Trello、Tableau、MS Project}
\cvline{溝通與領導}{專家 \hfill \textbf{關鍵字}：跨部門協調、利害關係人管理、團隊激勵、簡報技巧、衝突解決}
\cvline{技術知識}{中級 \hfill \textbf{關鍵字}：SQL、API、CI/CD、雲端服務、資料分析}
\section{榮譽}

\cventry{2023年12月}
        {智邦科技股份有限公司}
        {年度最佳專案經理}
        {}
        {}
        {表彰在企業級網通軟體專案中展現卓越領導力、風險管理與跨部門協調能力，並成功將交付週期縮短 33\%。}

\cventry{2020年11月}
        {趨勢科技股份有限公司}
        {專案管理卓越獎}
        {}
        {}
        {肯定在雲端安全產品全球發布專案中的傑出表現，達成 100\% 里程碑並維持客戶滿意度 95\% 以上。}
\section{證書}

\cventry{2018年6月}
        {Project Management Institute}
        {Project Management Professional (PMP)}
        {\url{https://www.pmi.org/certifications/project-management-pmp}}
        {}
        {}

\cventry{2019年9月}
        {Scrum Alliance}
        {Certified ScrumMaster (CSM)}
        {\url{https://www.scrumalliance.org/certifications/csm}}
        {}
        {}

\cventry{2021年4月}
        {Project Management Institute}
        {PMI Agile Certified Practitioner (PMI-ACP)}
        {\url{https://www.pmi.org/certifications/agile-acp}}
        {}
        {}
\section{出版刊物}

\cventry{2016年6月}
        {敏捷式專案管理在台灣軟體產業之應用}
        {國立臺灣大學資訊管理學系碩士論文}
        {\url{https://hdl.handle.net/11296/example}}
        {}
        {\begin{itemize}
\item 探討敏捷式方法在台灣軟體公司的導入現況與挑戰。
\item 訪談 12 家企業、回收 180 份問卷，提出混合式管理框架。
\item 研究結果獲得多家受訪企業採用於內部流程改善。
\end{itemize}}
\section{項目}

\cventry{2022年1月–2023年8月}
        {整合需求管理、測試追蹤與版本發布的內部專案管理平台}
        {企業級網通軟體交付平台}
        {\url{https://www.accton.com.tw}}
        {}
        {\begin{itemize}
\item 主導平台導入，涵蓋 5 個產品線、超過 40 位使用者。
\item 建立自動化儀表板，即時呈現專案健康度與資源負載。
\item 縮短 30\% 專案狀態彙報時間，並提升資料透明度。
\end{itemize}
\textbf{關鍵字}：專案管理平台、數據視覺化、流程自動化}

\cventry{2019年3月–2020年12月}
        {面向全球企業客戶的雲端安全解決方案發布計畫}
        {雲端安全產品全球發布專案}
        {\url{https://www.trendmicro.com}}
        {}
        {\begin{itemize}
\item 協調台灣、美國、日本與印度團隊，完成 6 個主要版本發布。
\item 建立跨時區溝通機制，將決策週期縮短 20\%。
\item 專案達成率 100\%，客戶滿意度維持 95\% 以上。
\end{itemize}
\textbf{關鍵字}：全球發布、跨時區協作、雲端安全}
\section{興趣愛好}

\cvline{戶外活動}{登山、馬拉松、單車環島}
\cvline{閱讀與學習}{商業策略、專案管理、人工智慧}
\cvline{社群參與}{敏捷社群、專案管理協會、技術講座}
\section{志願服務}

\cventry{2019年1月–2023年12月}
        {志工講師}
        {台灣專案管理學會}
        {\url{https://www.pmi.org.tw}}
        {}
        {\begin{itemize}
\item 擔任敏捷與專案管理課程志工講師，累計授課 30 場、學員超過 600 人。
\item 協助籌辦年度專案管理論壇，負責議程規劃與講者邀請。
\item 參與偏鄉資訊教育志工活動，指導高中生專題製作與簡報技巧。
\end{itemize}}
    
\end{document}